\chapter{Figuras, Tabelas, Equações e Outros Objetos}

\section{Figura individual}

\begin{lstlisting}
\undffigura[0.75\textwidth]
  {figuras/exemplo.jpg}
  {Legenda da figura}
  {rotulo}
  [Autor, ano]
\end{lstlisting}

A largura e a fonte são opcionais. Sem indicação, o comando usa \texttt{0.8\string\textwidth} e \texttt{elaboração própria}. Referencie com \texttt{\string\figuraref\{rotulo\}}.

\section{Duas figuras independentes lado a lado}

\begin{lstlisting}
\undffigurasladoalado
  {figuras/antes.jpg}{Registro anterior}{antes}[elaboração própria]
  {figuras/depois.jpg}{Registro posterior}{depois}[elaboração própria]
\end{lstlisting}

O comando cria duas figuras independentes, cada uma com número, legenda, rótulo e fonte. Não há legenda coletiva nem subdivisões ``a'' e ``b''. A largura opcional inicial controla cada coluna.

\section{Tabela}

\begin{lstlisting}
\begin{undftabela}{Legenda da tabela}{rotulo}[Fonte]
  \begin{tabular}{ll}
    Item & Valor \\
  \end{tabular}
\end{undftabela}
\end{lstlisting}

Use \texttt{\string\tabelaref\{rotulo\}} para a chamada no texto. Tabelas extensas que atravessam páginas devem ser construídas com um ambiente multipágina apropriado, como \texttt{longtable}, pois \texttt{undftabela} é um flutuante de uma página.

\section{Gráficos e quadros}

Na versão 0.1.0, um gráfico fornecido como imagem é inserido por \texttt{\string\undffigura}; a saída o identifica na categoria geral ``Figura''. Um gráfico produzido diretamente em LaTeX pode ser colocado em um ambiente \texttt{figure}, com \texttt{\string\caption}, rótulo iniciado por \texttt{fig:} e \texttt{\string\undffonte}.

Não há ambiente público dedicado a quadro na versão 0.1.0. Não use \texttt{undftabela} apenas para obter bordas quando o objeto for semanticamente um quadro, pois a identificação sairia como ``Tabela''. Até que uma interface própria seja incorporada e testada, quadros exigem implementação manual controlada pelo mantenedor. Essa limitação está registrada para uma versão posterior e não deve ser ocultada por documentação de um comando inexistente.

\section{Equações}

\begin{lstlisting}
\equacao[energia]{E = mc^2}
\equacaosemnumero{x + y = z}
\end{lstlisting}

Somente a equação numerada e rotulada pode ser referenciada:

\begin{lstlisting}
\equacaoref{energia}
\end{lstlisting}

\section{Fluxograma}

\begin{lstlisting}
\begin{undffluxograma}{Fluxo de análise}{analise}
  \centering
  \begin{tikzpicture}
    % nós e conexões
  \end{tikzpicture}
\end{undffluxograma}
\end{lstlisting}

Até três ocorrências, o recurso integra a numeração e a lista de figuras. Acima desse total, passa a usar numeração e lista próprias. Após adicionar ou remover fluxogramas, compile novamente até os rótulos e as listas estabilizarem.

\section{Pseudocódigo}

\begin{lstlisting}
\begin{undfpseudocodigo}{MEDIA(V)}{media}{Calcular a média}
Entrada: lista V
Saida: média dos valores
...
\end{undfpseudocodigo}
\end{lstlisting}

O ambiente comum tenta iniciar o bloco apenas quando há espaço suficiente para cerca de meia página. Para um algoritmo que precise continuar na página seguinte, use \texttt{undfpseudocodigolongo}; ele permite a quebra interna do bloco. Referencie ambos com \texttt{\string\pseudocodigoref\{rotulo\}}.

\section{Apêndices, anexos e PDFs externos}

\begin{lstlisting}
\undfapendice{Título}{rotulo}
\undfanexo{Título}{rotulo}
\undfapendicepdf{Título}{arquivo.pdf}
\undfanexarpdf{Título}{arquivo.pdf}
\end{lstlisting}

Para inclusão sem criar apêndice ou anexo, use \texttt{\string\undfincluirpdf\{arquivo.pdf\}}. Os comandos de referência são \texttt{\string\apendiceref} e \texttt{\string\anexoref}.

\section{Imagens e paginação}

Os comandos de figura verificam o tamanho do arquivo e emitem avisos para imagens pesadas. Uma figura fixada que não caiba no espaço restante avança inteira para a próxima página, podendo deixar área em branco. Isso preserva a unidade do objeto. Ajuste a largura, reposicione a chamada no texto ou use uma imagem mais compacta; não force sobreposição com texto ou margens.
